% Quiz 06 --- Find a graph whose four centers are four different nodes.
% Printed, worked in class, 20 minutes. One question, 10 points. Page 1 is the
% question and a box to draw in; page 2 is a work page with a table.
%
% The task is the Brandes--Hildenbrand (2014, Network Science 2(3), 416--418)
% "distinct singleton centers" problem: a connected graph in which betweenness,
% closeness, degree and eigenvector centrality each have exactly one most
% central node, and the four nodes are different. The smallest such graphs have
% 10 nodes; the node limit on this sheet is 12 (see README.md for why).
%
% The sheet does not say "use code", and does not forbid it. Do not add either.
% It also does not name the graph class (the paper calls it BCDE): the four
% measures are named in full, so there is no letter to confuse with a node.
%
% go.skojaku.com does not point at this file; the answers are in solutions.tex.
%
% Compile with: xelatex -interaction=nonstopmode quiz06.tex   (run twice)
\documentclass[a4paper, 11pt]{extarticle}
% Shared preamble for Quiz 06 (quiz06.tex and solutions.tex).
% Compile with: xelatex -interaction=nonstopmode <file>.tex   (run twice)
%
% Same house look as the pen-and-paper sheets in
% lecture-note/*/pen-and-paper/: Charter, A4, generous leading, room to draw.

\usepackage[top=0.7in, bottom=0.6in, left=0.8in, right=0.8in]{geometry}
\usepackage{amsmath}
\usepackage{amssymb}
\usepackage{array}
\usepackage{graphicx}
\usepackage[hidelinks]{hyperref}
\usepackage{tcolorbox}
\usepackage{enumitem}
\usepackage{etoolbox}
\usepackage{tikz}
\usepackage{fontspec}
\usetikzlibrary{calc,positioning,arrows.meta}
\definecolor{pltblue}{HTML}{1F77B4}
\definecolor{rulegrey}{HTML}{6A6D75}

% Body face. Charter, and the four ways to get it. XCharter *is* Charter and
% ships with a full TeX Live, so it is asked for first and by file: a CI runner
% answers yes to \IfFontExistsTF{Charter}, then cannot typeset it.
\IfFontExistsTF{XCharter-Roman.otf}{
  \setmainfont{XCharter}[
    Extension      = .otf,
    UprightFont    = *-Roman,
    BoldFont       = *-Bold,
    ItalicFont     = *-Italic,
    BoldItalicFont = *-BoldItalic]
}{
  \IfFontExistsTF{Charter}{\setmainfont{Charter}}{
    \IfFontExistsTF{Iowan Old Style}{\setmainfont{Iowan Old Style}}{
      \IfFontExistsTF{Georgia}{\setmainfont{Georgia}}{}}}}

\setlength{\parindent}{0pt}
\setlength{\parskip}{0.4em}
\setlength{\emergencystretch}{1.2em}
\raggedbottom
\pagestyle{empty}

% Quiz header. #1 = title, #2 = the line under it, #3 = term.
\newcommand{\quizhead}[3]{%
  {\footnotesize\textsc{Advanced Network Science} \hfill \textsc{#3}}\par
  \vspace{-0.2em}\rule{\textwidth}{1pt}\par
  \vspace{0.3em}
  {\LARGE\bfseries #1}\par
  \vspace{0.1em}
  {\small #2}\par
  \vspace{0.7em}
  {\small Name \underline{\hspace{6cm}} \hfill Date \underline{\hspace{3.2cm}}}\par
  \vspace{0.7em}}

% A question heading with its point value.
\newcommand{\question}[3]{%
  \vspace{0.5em}
  {\large\bfseries Question #1 \quad\normalsize\mdseries (#2 points)}\par
  \vspace{-0.3em}{\color{rulegrey}\rule{\textwidth}{0.6pt}}\par
  \vspace{0.1em}
  {\bfseries #3}\par
  \vspace{0.2em}}

% An empty box to draw or write in. #1 = height, #2 = caption under the rule.
\newcommand{\drawbox}[2]{%
  \begin{tcolorbox}[colback=white, colframe=rulegrey, boxrule=0.5pt,
                    sharp corners, height=#1, valign=top,
                    boxsep=2pt, left=4pt, right=4pt, top=3pt, bottom=3pt]
  {\footnotesize\color{rulegrey} #2}
  \end{tcolorbox}}

% Ruled writing lines, #1 of them.
\newcommand{\writelines}[1]{%
  \vspace{0.3em}
  \foreach \i in {1,...,#1}{%
    {\color{rulegrey}\rule{\textwidth}{0.4pt}}\par\vspace{1.5em}}}


% Title and points share one line, as in Quiz 5.
\renewcommand{\question}[3]{%
  \vspace{0.4em}
  {\large\bfseries Question #1 \quad #3\quad\normalsize\mdseries (#2 points)}\par
  \vspace{-0.4em}{\color{rulegrey}\rule{\textwidth}{0.6pt}}\par}

\begin{document}

\quizhead{Quiz 6 \; --- \; Four centers}
        {20 minutes \quad$\cdot$\quad 10 points}
        {Fall 2026}

% ---------------------------------------------------------------------------
\question{1}{10}{Find a graph with four different centers.}

\begin{tcolorbox}[colback=pltblue!4, colframe=pltblue!35, boxrule=0.5pt,
                  sharp corners, boxsep=2pt, left=6pt, right=6pt,
                  top=4pt, bottom=4pt]
{\small
Take a connected, undirected graph with $N$ nodes and no self-loops or
multiple edges. Write $d_{ij}$ for the shortest-path distance from $i$ to $j$,
and $\sigma_{st}$ for the number of shortest paths from $s$ to $t$, of which
$\sigma_{st}(i)$ pass through $i$. For a node $i$:
\begin{center}
\renewcommand{\arraystretch}{1.35}
\begin{tabular}{@{}ll@{}}
\textbf{degree} & $k_i$, the number of neighbors of $i$ \\
\textbf{closeness} & $\displaystyle \frac{N-1}{\sum_{j} d_{ij}}$ \\
\textbf{betweenness} & $\displaystyle \sum_{s<t,\; s,t \neq i} \frac{\sigma_{st}(i)}{\sigma_{st}}$ \\
\textbf{eigenvector} & $x_i$, where $A x = \lambda_{\max} x$ and every $x_i > 0$ \\
\end{tabular}
\end{center}
The \emph{center} under a measure is the node with the largest value. The
center is \emph{unique} when no other node ties with it.
}
\end{tcolorbox}

\vspace{0.3em}
Draw a graph with \textbf{at most 12 nodes} such that
\begin{itemize}[leftmargin=1.4em, itemsep=0.1em, topsep=0.2em]
  \item under each of the four measures, the center is unique, and
  \item the four centers are four different nodes.
\end{itemize}

(a) [2 pts] Draw the graph below. Number the nodes $1, 2, \dots$ The graph must
be connected, with at most 12 nodes.

(b) [8 pts] For each of the four measures, name the center and show that it
beats every other node. Use the table on the back if it helps.

\vspace{0.2em}
\drawbox{7.4cm}{Your graph}

\vfill

{\footnotesize\color{rulegrey}
Done early? (not graded) Brandes and Hildenbrand (2014) searched every
connected graph with fewer than 10 nodes by computer and found no graph with
four different unique centers. How few nodes can you manage?}

% ===========================================================================
\newpage

{\large\bfseries Work page}\par
\vspace{-0.4em}{\color{rulegrey}\rule{\textwidth}{0.6pt}}\par
\vspace{0.4em}

{\small Fill in as much as you need. Then circle each center.}

\vspace{0.5em}
\renewcommand{\arraystretch}{1.9}
\begin{center}
\begin{tabular}{|c|c|c|c|c|}
\hline
\textbf{node} & \textbf{degree} & $\sum_j d_{ij}$ \textbf{(for closeness)} &
\textbf{betweenness} & \textbf{eigenvector} \\
\hline
1  & & & & \\ \hline
2  & & & & \\ \hline
3  & & & & \\ \hline
4  & & & & \\ \hline
5  & & & & \\ \hline
6  & & & & \\ \hline
7  & & & & \\ \hline
8  & & & & \\ \hline
9  & & & & \\ \hline
10 & & & & \\ \hline
11 & & & & \\ \hline
12 & & & & \\ \hline
\end{tabular}
\end{center}

\vspace{0.6em}
{\small
For eigenvector centrality, you do not need $\lambda_{\max}$ exactly. Any
values of $x$ that you can justify will do, as long as they show who is
largest.}

\end{document}
